\chapter{Einleitung}
\label{ch:einleitung}

\section{Einordnung in das Gesamtprojekt}
\label{sec:einordnung}

Das Projekt \enquote{Dynamisches Greifen und automatisierte Sichtprüfung mit einem
kollaborativen Roboter} ist in vier Module gegliedert. Modul~1 erkennt das Bauteil mit
einer Draufsichtkamera und bestimmt seine Lage, Modul~2 steuert den Roboter und greift das
Bauteil auf klassischem Weg, und Modul~3 prüft das abgelegte Bauteil optisch. Modul~4, das
Thema dieses Dokuments, untersucht einen anderen Weg zum Greifen: Statt einer programmierten
Bahn zu einer gemessenen Greifpose führt eine gelernte Steuerungsstrategie (Policy) den Arm.
Sie erhält in jedem Takt Kamerabilder und den Zustand des Roboters und gibt daraus direkt die
nächsten Sollwerte für Gelenke und Greifer aus.

Modul~4 ist eine Alternative zum Greifen von Modul~2 und für den automatischen Ablauf nicht
erforderlich. Im Ablauf übernimmt es nur den Abschnitt vom Start bis zur Übergabe: Die Policy
startet in einer festen Ausgangsstellung, sucht das Bauteil in den Bildern, greift es und
fährt es an einen Übergabepunkt über der Prüfvorrichtung. Dort endet ihre Aufgabe mit dem
Bauteil im Greifer. Das Ablegen auf der Prüfvorrichtung und alle weiteren Bewegungen führt
wieder Modul~2 aus (\cref{ch:integration}).

\section{Zielsetzung}
\label{sec:ziel}

Ziel von Modul~4 ist eine vollständige, lokal betriebene Kette für Imitation Learning am
Roboter Neura LARA~5:
\begin{enumerate}
  \item automatisierte Aufnahme von Demonstrationen mit gezielt eingespeistem Rauschen,
  \item Training einer Diffusion Policy~\cite{chi2024} auf diesen Daten,
  \item Ausführung der Policy in Echtzeit mit einer eigenen Sicherheitsschicht und
  \item Einbindung in den Ablauf und die Bedienoberfläche des Gesamtprojekts.
\end{enumerate}
Bewertet wird die Kette danach, wie zuverlässig die Policy das Bauteil an unterschiedlichen
Lagen greift und übergibt (\cref{ch:versuche}).

\section{Aufbau der Anlage}
\label{sec:aufbau}

\Cref{tab:aufbau} fasst die Komponenten zusammen, mit denen Modul~4 arbeitet. Roboter und
Greifer teilt es mit Modul~2, die Draufsichtkamera mit Modul~1.

\begin{table}[htbp]
  \centering
  \caption{Komponenten von Modul~4}
  \label{tab:aufbau}
  \begin{tabularx}{\textwidth}{@{}l X@{}}
    \toprule
    Komponente & Ausführung \\
    \midrule
    Roboter & Neura LARA~5, sechs Achsen, Steuerung LARA-5.0.8, Python-Schnittstelle
              NeuraPy über TCP/IP~\cite{neura2025neurapy} \\
    Greifer & Zimmer LWR50L-03, Parallelbackengreifer, über NeuraPy nur öffnen und
              schließen~\cite{zimmer_lwr50l} \\
    Handgelenkkamera & Daheng VEN-161-61U3C, $1440 \times 1080$~Pixel, Global Shutter,
                       Objektiv \SI{3.6}{\milli\metre}, F8.0, kaum Verzeichnung \\
    Szenenkamera & Draufsichtkamera von Modul~1 \\
    Bauteil & $\num{39.7} \times \SI{14.9}{\milli\metre}$ (Grundfläche) \\
    Rechner & Training: NVIDIA RTX 5070~Ti, \SI{16}{\giga\byte}; Betrieb: Laptop mit
              NVIDIA-\acs{gpu} \\
    \bottomrule
  \end{tabularx}
\end{table}

\todo[inline]{Foto des Aufbaus mit Beschriftung (Roboter, Kameras, Bauteil, Prüfvorrichtung,
Kamerastange).}

\section{Aufbau dieses Dokuments}
\label{sec:aufbau-dokument}

\Cref{ch:grundlagen} stellt den klassischen Ansatz und Behavior Cloning gegenüber und
erklärt die mathematischen Grundlagen der verwendeten Verfahren, einschließlich der bekannten
Schwierigkeiten des Lernens aus Demonstrationen. \Cref{ch:anforderungen} leitet daraus die
Anforderungen an Modul~4 ab. Die \cref{ch:architektur,ch:datenaufnahme,ch:training,%
ch:ausfuehrung,ch:sicherheit,ch:integration} beschreiben die Umsetzung. Jeder Abschnitt folgt
dabei demselben Muster: Idee, Umsetzung, aufgetretene Probleme und ihre Lösung, und die
Begründung der jeweiligen Designentscheidung. \Cref{ch:versuche} fasst die Versuche
zusammen, \cref{ch:diskussion} bewertet sie, und \cref{ch:fazit} zieht ein Fazit.

Die Messberichte und Einzeldokumente im Ordner \code{code/m4_behavior_cloning/} bleiben als
eigenständige Unterlagen bestehen. Dieses Dokument fasst ihren Inhalt zusammen und verweist
für Rohdaten und Einzelheiten auf sie (\cref{app:berichte}).
